% GENERATED FILE. Edit canonical lessons, then run npm run generate:books.
% canonical-source-hash: fnv1a64:3cc4780203bc8645
% canonical-lessons: RU-C363-peregovory, RU-C363-kontrakt, RU-C363-zakazchik, RU-C363-konkurent, RU-C363-sotrudnichat

\chapter{Negotiations and contracts}
\label{ch:ru-negotiations-and-contracts}
\hlchaptermodality{363}

\begin{chapteropening}
\textbf{By the end of this chapter:} \emph{I can talk about negotiations, a contract, a client and a competitor, and say that two firms cooperate.}

Complete the last lesson of chapter 363: I can talk about negotiations, a contract, a client and a competitor, and say that two firms cooperate.
\end{chapteropening}

\section[\texorpdfstring{\ru{переговоры} (peregovóry)}{переговоры (peregovóry)}]{\ru{переговоры} (peregovóry) — negotiations, talks}

\label{lesson:RU-C363-peregovory}



\begin{quote}
Before the new one: say the Russian for a specialist, then the Russian for to dismiss.
\end{quote}

\subsection*{You'll want to know: \ru{переговоры}}
\textbf{\ru{переговоры}} (\emph{peregovóry}) — \textquotedblleft{}negotiations, talks\textquotedblright{}.

\begin{quote}
\ru{Переговоры} \ru{шли} \ru{всю} \ru{ночь}.
\end{quote}

\begin{quote}
\emph{Peregovóry shli vsyu noch.}
\end{quote}

\begin{quote}
The talks went on all night.
\end{quote}

Plural only. \emph{\ru{Вести} \ru{переговоры}} — to hold talks.

\begin{cousinweb}
\textbf{\ru{переговоры}} is built from \emph{\ru{пере}-} \textquotedblleft{}across\textquotedblright{} and \emph{\ru{говор}} \textquotedblleft{}speech\textquotedblright{}, the root of \emph{\ru{говорить}}: talk that goes back and forth across a table.
\end{cousinweb}

\begin{grammarlens}[title={What you've built}]
One more word for telling, explaining and describing at B1.
\end{grammarlens}

\subsection*{Your turn}
\begin{itemize}
\raggedright
  \item \emph{Say it:} \emph{peregovóry}
  \item \emph{Say it:} \emph{peregovóry}, once more
  \item \emph{Say it:} a sentence of your own with \emph{peregovóry}
  \item \emph{Recall:} say the Russian for a specialist, then the Russian for to dismiss, then say \emph{peregovóry} again
\end{itemize}

\begin{tcolorbox}[breakable,colback=teal!4,colframe=teal!35!black,title={Before you move on}]
What does \ru{переговоры} mean? (\textquotedblleft{}Negotiations, talks\textquotedblright{}.) Say it once more.
\end{tcolorbox}
\section[\texorpdfstring{\ru{контракт} (kontrákt)}{контракт (kontrákt)}]{\ru{контракт} (kontrákt) — a contract}

\label{lesson:RU-C363-kontrakt}



\begin{quote}
Before the new one: say the Russian for to dismiss, then the Russian for negotiations.
\end{quote}

\subsection*{You'll want to know: \ru{контракт}}
\textbf{\ru{контракт}} (\emph{kontrákt}) — \textquotedblleft{}a contract\textquotedblright{}.

\begin{quote}
\ru{Он} \ru{подписал} \ru{контракт} \ru{на} \ru{три} \ru{года}.
\end{quote}

\begin{quote}
\emph{On podpisál kontrákt na tri góda.}
\end{quote}

\begin{quote}
He signed a three-year contract.
\end{quote}

Masculine. \emph{\ru{Работать} \ru{по} \ru{контракту}} — to work on a contract. Close to \emph{\ru{договор}}.

\begin{cousinweb}
\textbf{\ru{контракт}} was borrowed from German \emph{Kontrakt}, from Latin \emph{contractus} (EtymDB) \textquotedblleft{}a drawing together, an agreement\textquotedblright{}, from \emph{contrahere} \textquotedblleft{}to draw together\textquotedblright{} (Lewis \& Short).
\end{cousinweb}

\begin{grammarlens}[title={What you've built}]
One more word for telling, explaining and describing at B1.
\end{grammarlens}

\subsection*{Your turn}
\begin{itemize}
\raggedright
  \item \emph{Say it:} \emph{kontrákt}
  \item \emph{Say it:} \emph{kontrákt}, once more
  \item \emph{Say it:} a sentence of your own with \emph{kontrákt}
  \item \emph{Recall:} say the Russian for to dismiss, then the Russian for negotiations, then say \emph{kontrákt} again
\end{itemize}

\begin{tcolorbox}[breakable,colback=teal!4,colframe=teal!35!black,title={Before you move on}]
What does \ru{контракт} mean? (\textquotedblleft{}A contract\textquotedblright{}.) Say it once more.
\end{tcolorbox}
\section[\texorpdfstring{\ru{заказчик} (zakázchik)}{заказчик (zakázchik)}]{\ru{заказчик} (zakázchik) — a client (who orders work)}

\label{lesson:RU-C363-zakazchik}



\begin{quote}
Before the new one: say the Russian for negotiations, then the Russian for a contract.
\end{quote}

\subsection*{You'll want to know: \ru{заказчик}}
\textbf{\ru{заказчик}} (\emph{zakázchik}) — \textquotedblleft{}a client (who orders work)\textquotedblright{}.

\begin{quote}
\ru{Заказчик} \ru{недоволен} \ru{нашей} \ru{работой}.
\end{quote}

\begin{quote}
\emph{Zakázchik nedovólen náshey rabótoy.}
\end{quote}

\begin{quote}
The client is unhappy with our work.
\end{quote}

Masculine. The one who orders a job or a product, from \emph{\ru{заказывать}}.

\begin{cousinweb}
\textbf{\ru{заказчик}} \textquotedblleft{}buyer, client\textquotedblright{} (EtymDB) is built on \emph{\ru{заказывать}} \textquotedblleft{}to order\textquotedblright{}, whose root is \emph{\ru{казать}} \textquotedblleft{}to show, to say\textquotedblright{}, Proto-Slavic \emph{*kazati} (EtymDB).
\end{cousinweb}

\begin{grammarlens}[title={What you've built}]
One more word for telling, explaining and describing at B1.
\end{grammarlens}

\subsection*{Your turn}
\begin{itemize}
\raggedright
  \item \emph{Say it:} \emph{zakázchik}
  \item \emph{Say it:} \emph{zakázchik}, once more
  \item \emph{Say it:} a sentence of your own with \emph{zakázchik}
  \item \emph{Recall:} say the Russian for negotiations, then the Russian for a contract, then say \emph{zakázchik} again
\end{itemize}

\begin{tcolorbox}[breakable,colback=teal!4,colframe=teal!35!black,title={Before you move on}]
What does \ru{заказчик} mean? (\textquotedblleft{}A client (who orders work)\textquotedblright{}.) Say it once more.
\end{tcolorbox}
\section[\texorpdfstring{\ru{конкурент} (konkurént)}{конкурент (konkurént)}]{\ru{конкурент} (konkurént) — a competitor}

\label{lesson:RU-C363-konkurent}



\begin{quote}
Before the new one: say the Russian for a contract, then the Russian for a client.
\end{quote}

\subsection*{You'll want to know: \ru{конкурент}}
\textbf{\ru{конкурент}} (\emph{konkurént}) — \textquotedblleft{}a competitor\textquotedblright{}.

\begin{quote}
\ru{Наш} \ru{главный} \ru{конкурент} \ru{снизил} \ru{цены}.
\end{quote}

\begin{quote}
\emph{Nash glávnyy konkurént snízil tsény.}
\end{quote}

\begin{quote}
Our main competitor has cut its prices.
\end{quote}

Masculine. \emph{\ru{Конкуренция}} — competition.

\begin{cousinweb}
\textbf{\ru{конкурент}} is usually traced, through German, to Latin \emph{concurrēns}, the participle of \emph{concurrere} \textquotedblleft{}to run together, to clash\textquotedblright{} (Lewis \& Short), built on \emph{currere} \textquotedblleft{}to run\textquotedblright{}.
\end{cousinweb}

\begin{grammarlens}[title={What you've built}]
One more word for telling, explaining and describing at B1.
\end{grammarlens}

\subsection*{Your turn}
\begin{itemize}
\raggedright
  \item \emph{Say it:} \emph{konkurént}
  \item \emph{Say it:} \emph{konkurént}, once more
  \item \emph{Say it:} a sentence of your own with \emph{konkurént}
  \item \emph{Recall:} say the Russian for a contract, then the Russian for a client, then say \emph{konkurént} again
\end{itemize}

\begin{tcolorbox}[breakable,colback=teal!4,colframe=teal!35!black,title={Before you move on}]
What does \ru{конкурент} mean? (\textquotedblleft{}A competitor\textquotedblright{}.) Say it once more.
\end{tcolorbox}
\section[\texorpdfstring{\ru{сотрудничать} (sotrúdnichat)}{сотрудничать (sotrúdnichat)}]{\ru{сотрудничать} (sotrúdnichat) — to cooperate, to work together}

\label{lesson:RU-C363-sotrudnichat}



\begin{quote}
Before the new one: say the Russian for a client, then the Russian for a competitor.
\end{quote}

\subsection*{You'll want to know: \ru{сотрудничать}}
\textbf{\ru{сотрудничать}} (\emph{sotrúdnichat}) — \textquotedblleft{}to cooperate, to work together\textquotedblright{}.

\begin{quote}
\ru{Наша} \ru{фирма} \ru{сотрудничает} \ru{с} \ru{немецкой} \ru{компанией}.
\end{quote}

\begin{quote}
\emph{Násha fírma sotrúdnichayet s nemétskoy kompániyey.}
\end{quote}

\begin{quote}
Our firm works together with a German company.
\end{quote}

Imperfective only. \emph{\ru{Сотрудничать} \ru{с}} + instrumental.

\begin{cousinweb}
\textbf{\ru{сотрудничать}} is built on \emph{\ru{сотрудник}}, from \emph{\ru{со}-} \textquotedblleft{}together\textquotedblright{} and \emph{\ru{труд}} \textquotedblleft{}work\textquotedblright{}: to do work together.
\end{cousinweb}

\begin{grammarlens}[title={What you've built}]
One more word for telling, explaining and describing at B1.
\end{grammarlens}

\subsection*{Your turn}
\begin{itemize}
\raggedright
  \item \emph{Say it:} \emph{sotrúdnichat}
  \item \emph{Say it:} \emph{sotrúdnichat}, once more
  \item \emph{Say it:} a sentence of your own with \emph{sotrúdnichat}
  \item \emph{Recall:} say the Russian for a client, then the Russian for a competitor, then say \emph{sotrúdnichat} again
\end{itemize}

\begin{tcolorbox}[breakable,colback=teal!4,colframe=teal!35!black,title={Before you move on}]
What does \ru{сотрудничать} mean? (\textquotedblleft{}To cooperate, to work together\textquotedblright{}.) Say it once more.
\end{tcolorbox}
